\section{项目背景与产品定位}

\subsection{自由行规划中的真实痛点}

中文自由行用户在筹备一趟三至七天的城市旅行时，往往要在四类割裂的工具间反复切换：浏览器收藏夹与 PDF 里攒着零散攻略，地图软件里逐段量算通勤时长，场馆公众号或官网里翻找淡旺季开放时间与预约规则，最后再在备忘录里手工拼出按天时间表。近年来，通用大语言模型虽然能够一键输出流畅的行程文案，但一旦拿来真正落地，四类结构性缺陷便迅速暴露：

\begin{itemize}
  \item \textbf{时效事实与凭空编造混为一谈}：模型容易把过期的淡旺季时间、第三方攻略中的模糊报价，甚至凭记忆拼凑的闭馆日直接写成笃定结论，用户既无法追溯来源，也难以分辨哪些数字是已核实的官方标准、哪些只是粗略估算。
  \item \textbf{地理与时间可行性缺乏硬约束}：行程里常见跨城折返、步行数公里却只留十余分钟衔接、活动结束晚于场馆停止入场时间等问题，读来顺畅却无法执行。
  \item \textbf{微调成本高，局部修改引发全局重写}：用户只想删掉一个景点、调换两项顺序或把某项活动移到第二天，重新对话往往导致其余早已满意的日期被整篇改写，预算也无法保持口径一致。
  \item \textbf{个人资料与长期习惯难以克制地融入}：用户手头的私人笔记与旅行偏好要么每次手动复制粘贴，要么被模型过度发挥——不仅把长期习惯凌驾于本次明确要求之上，还在页面中输出大量冗余的“因为您偏好……所以推荐……”式解释文本。
\end{itemize}

\subsection{产品定位与核心边界}

针对上述问题，本项目（“拾程”）定位为一款\textbf{单人本机运行、数据本地持久化、事实可核对且支持局部编辑}的中文自由行规划 Agent。在需求收敛阶段，我们以 24 条显式产品规则（R01--R24）划定了清晰的能力边界，避免范围蔓延：

\begin{enumerate}
  \item \textbf{场景聚焦}：面向目的地与起止日期已确定的 \textbf{3--7 天中国大陆单城休闲旅行}，覆盖景点、餐饮与当地交通；航班、高铁与住宿作为用户提供的已知条件参与时间与位置约束，首版不代为搜索酒店、机票或执行预订。
  \item \textbf{本机优先与单人安装}：每份安装供个人在自有电脑上通过本机浏览器使用，所有资料、偏好、行程草案与历史版本均保存在本机数据目录或持久化卷中，无需依赖云端账户体系（ADR-0002、ADR-0003）。
  \item \textbf{诚实核验，未知显式呈现}：营业时间、票价与预约优先采信官方来源，时效事实必须关联来源 URL 与查询时间；无法可靠核实的内容明确标记为“待核对”或“估算”，绝不将未知包装为已满足条件（R10--R12）。
  \item \textbf{可编辑与定点重排}：用户可在草案上执行增删活动、替换地点、同日排序与明确指定的跨日移动；自动重排与事实/路线刷新严格限定在受影响日期，其余日期逐字段保持原值，全团全程预算同步重算（R21）。
  \item \textbf{清晰的个性化优先级}：遵循“\textbf{本次明确要求 $>$ 长期偏好档案 $>$ 个人参考资料}”的层级；行程中的临时调整不会悄悄改写长期偏好，只有经用户主动确认保存后才会影响后续新行程，且页面绝不渲染多余的偏好匹配理由（R01--R04，ADR-0001）。
\end{enumerate}

\begin{summarybox}{首版质量口径与简化原则}
  本项目不以“单次生成耗时”或“API 调用费用”作为验收门槛，而是将质量标准锚定在三个可感知的目标上：\textbf{明确约束严格遵守、事实依据清晰可核对、用户只需少量修改便愿意采用}。工程实现遵循“简化原则”：以单个应用容器交付、以 SQLite + 本地中文分词替代重型向量数据库、以停服离线校验归档替代复杂在线容灾，把研发精力集中在规划质量与核验可靠性上。
\end{summarybox}

\subsection{领域核心概念对齐}

为避免前后端与规划器在语义上出现偏差，系统在术语表与数据契约中严格区分以下核心概念：

\begin{table}[H]
  \centering
  \caption{核心领域概念与业务边界}
  \renewcommand{\arraystretch}{1.28}
  \begin{tabularx}{\textwidth}{@{} >{\raggedright\arraybackslash\sffamily\bfseries}p{3.7cm} Y @{}}
    \toprule
    概念 & 定义与约束规则 \\
    \midrule
    行程要求\\{\footnotesize\texttt{TripRequest}} & 包含目的地、3--7 天日期、同行人数、可选全团预算（人民币分，覆盖门票、餐饮与当地交通，不含住宿与往返大交通）、自然语言补充说明、结构化抵离/预约时间及关联参考资料；每次修改生成新修订版。 \\
    旅行偏好\\{\footnotesize\texttt{PreferenceProfile}} & 单实例长期旅行习惯，以自然语言原文保存；仅在用户主动编辑保存后变更，带版本冲突保护。 \\
    个人参考资料\\{\footnotesize\texttt{KnowledgeSource}} & 用户主动导入的 PDF、文本、公开网页链接或收藏；按需关联到具体行程，仅将命中的必要片段发送给模型，绝不发送给地图或搜索引擎。 \\
    事实依据\\{\footnotesize\texttt{FactEvidence}} & 与具体地点、适用日期绑定的营业时间、票价、餐费或预约记录，明确区分 \texttt{supported}（有官方原文依据）、\texttt{estimated}（估算或非官方）、\texttt{conflict}（来源冲突）与 \texttt{unknown}（待核对）。 \\
    自由用餐时段\\{\footnotesize\texttt{flexible\_meal}} & 餐饮默认仅安排午餐/晚餐约略时段，以附近景点或街区作为交通锚点并标明餐费估算，不指定具体餐馆，也不套用锚点景点的营业时间与门票。 \\
    \bottomrule
  \end{tabularx}
\end{table}
